\chapter{Corrigés des exercices}\label{annexe:corriges}
\iterbanner{V3}

\section*{Chapitre~\ref{ch:question} — La question}
\textbf{Ex.~1.1.} \code{quadz} : un QUBO (après substitution $z=a_0+Ab$), 15 paramètres par grandeur, et il connaît l'ordre des niveaux puisque $z$ est proportionnel à l'indice. \code{fm} : un QUBO par construction, $1+10+10\times4=51$ paramètres, n'a aucune notion d'ordre. \code{bocs} : chaque tirage est un QUBO, $1+10+45=56$ paramètres, pas d'ordre non plus. \code{mlp} : pas un QUBO, 1\,249 paramètres, connaît l'ordre (son entrée est $z$).

\textbf{Ex.~1.2.} La boucle V2 suit exactement le schéma FMQA : apprendre un modèle quadratique en bits, le résoudre comme un QUBO, simuler le candidat, réapprendre. La différence tient à la forme des couplages. Dans le modèle en $z$, ils valent $2(A^\top MA)_{kl}$ : ils sont tous déterminés par les 15 coefficients du polynôme. Par exemple, deux bits $j$ et $l$ d'une même variable sont couplés par $2M_{ii}a_ja_l$, proportionnel à $2^{j+l}$ : la structure binaire est imposée. Dans une machine de factorisation, chaque couplage $\langle\mathbf v_k,\mathbf v_l\rangle$ est libre, à la contrainte de rang près.

\section*{Chapitre~\ref{ch:fm} — Machine de factorisation}
\textbf{Ex.~2.1.} $\hat y(\mathbf b)=1-b_1+0{,}5b_2-2b_3+0{,}5b_1b_2-b_1b_3$. Les huit valeurs ($b_1b_2b_3$) : $000\to1$ ; $100\to0$ ; $010\to1{,}5$ ; $110\to1$ ; $001\to-1$ ; $101\to-3$ ; $011\to-0{,}5$ ; $111\to-2$. Le minimum est $101$ ($-3$), vérifié dans le cours par la formule rapide.

\textbf{Ex.~2.2.} (a) En rang 1, $J_{ij}=v_iv_j$ (nombres réels), donc $J_{12}J_{13}J_{23}=(v_1v_2v_3)^2\ge0$, alors que $(-1)^3=-1$. (b) $\mathbf v_i=\sqrt2(\cos\theta_i,\sin\theta_i)$ avec $\theta_i=0^\circ,120^\circ,240^\circ$ : $\langle\mathbf v_i,\mathbf v_j\rangle=2\cos120^\circ=-1$. (c) Le rang limite les structures de couplages représentables, ici des « frustrations » où trois bits se repoussent deux à deux. Un rang trop petit empêche de représenter certaines physiques ; un rang trop grand ajoute des paramètres. On le règle par validation ; le code prend $k=4$.

\textbf{Ex.~2.3.} $3=110$ et $4=001$ diffèrent par leurs 3 bits ; $0=000$ et $7=111$ aussi ; $6=011$ et $7=111$ par un seul. Le niveau 4 est le seul où le bit $b_2$ est allumé seul. Sa prédiction vaut la constante plus le poids de $b_2$, et ce poids est appris sur les niveaux 5, 6 et 7, où $b_2$ est allumé \emph{avec} d'autres bits et leurs couplages. Sans profil au niveau 4, rien ne contraint la combinaison « $b_2$ seul » : c'est une extrapolation. Le modèle en $z$, lui, impose une parabole en $z$ qui passe nécessairement entre les niveaux 3 et 5 : il interpole.

\section*{Chapitre~\ref{ch:bocs} — Modèle bayésien}
\textbf{Ex.~3.1.} Le gradient de $\beta\|\Phi\boldsymbol\theta-\mathbf y\|^2+\boldsymbol\theta^\top A\boldsymbol\theta$ vaut $2\beta\Phi^\top(\Phi\boldsymbol\theta-\mathbf y)+2A\boldsymbol\theta$. Il s'annule pour $(A+\beta\Phi^\top\Phi)\boldsymbol\theta=\beta\Phi^\top\mathbf y$, soit $\boldsymbol\theta=\beta S\Phi^\top\mathbf y=\mathbf m$. En divisant par $\beta$ : c'est la solution ridge avec $\lambda=\alpha/\beta$ sur tous les coefficients sauf la constante.

\textbf{Ex.~3.2.} (a) $P(\theta<0)=\Phi(-0{,}3/0{,}2)=\Phi(-1{,}5)\approx0{,}07$. (b) $\Phi(-0{,}3)\approx0{,}38$. (c) En début de boucle, peu de données : la covariance $S$ est large, les tirages s'écartent souvent de la moyenne, et des régions jugées moyennes par $\mathbf m$ sont tirées comme bonnes, puis explorées. Quand les données s'accumulent, $S$ rétrécit et les tirages convergent vers $\mathbf m$ : on exploite.

\textbf{Ex.~3.3.} $\Phi$ a $N=30$ lignes, donc un rang au plus 30 : $\Phi^\top\Phi$ a au plus 30 valeurs propres non nulles. Les termes correspondant à $\lambda_k=0$ sont nuls, et chacun des autres est inférieur à 1 : $\gamma\le30$. Pour un modèle à 15 coefficients dont toutes les valeurs propres sont grandes devant $\alpha/\beta$, chaque terme vaut presque 1 : $\gamma\approx15$, tous les coefficients sont déterminés.

\section*{Chapitre~\ref{ch:mlp} — Réseau de neurones}
\textbf{Ex.~4.1.} $dh+h$ (première couche) $+\,h^2+h$ (deuxième) $+\,h+1$ (sortie) $=h^2+(d+3)h+1$. Pour $h=32$ : $1\,249$ ($d=4$) et $1\,441$ ($d=10$). Pour $d=10$, $h^2+13h+1=66$ donne $h\approx4$ (69 paramètres) : un réseau de 4 neurones par couche a autant de paramètres que le modèle quadratique en $z$.

\textbf{Ex.~4.2.} (a) Pour $f=b_0b_1b_2$, seule la chaîne $111$ contribue, avec le signe $(-1)^0=+1$ : le coefficient vaut 1. (b) En sommant avec les signes $(-1)^{3-|\mathbf b|}$ et la correspondance $x=b_0+2b_1+4b_2$, on trouve environ $-0{,}86$ : un terme cubique important, que le QUBO ne peut pas représenter. (c) $z_m$ est affine en bits, donc $z_m^2$ est quadratique : le coefficient cubique est nul (le calcul donne bien 0).

\section*{Chapitre~\ref{ch:precision} — Précision}
\textbf{Ex.~5.1.} $0{,}35\times10{,}4\approx3{,}6$ counts. Les profils voisins du front de S2 sont espacés de 0,1 à 3,9 counts (38,0 ; 40,4 ; 44,3 ; 44,4 ; 47,0 ; 50,9 ; 53,5 ; 55,8 ; 58,7) : l'erreur est du même ordre que l'écart entre voisins. Le modèle situe bien la région du front, mais il ne peut pas départager sûrement deux profils voisins. C'est pourquoi la boucle simule les candidats au lieu de croire le modèle.

\textbf{Ex.~5.2.} La corrélation de rang ne regarde que l'ordre des valeurs, pas leur écart. Un modèle qui prédit $\hat y=y+\sigma$ (la vraie valeur décalée d'un écart-type) classe parfaitement les profils (rang 1) mais se trompe de $\sigma$ partout : erreur relative 1. Pour choisir le meilleur profil, ce modèle serait parfait.

\section*{Chapitre~\ref{ch:contraintes} — Contraintes}
\textbf{Ex.~6.1.} Faisables entre $0{,}3$ et $1{,}1$ : les infaisables reçoivent $1{,}1+0{,}2\times0{,}8=1{,}26$. Si un nouveau faisable vaut $1{,}5$, le maximum et l'étendue changent : les infaisables valent maintenant $1{,}5+0{,}2\times1{,}2=1{,}74$. Les données des infaisables bougent d'une itération à l'autre ; le modèle apprend une cible qui se déplace.

\textbf{Ex.~6.2.} Problème V2 : 14 qubits contre 11, soit 3 qubits d'écart en moins, et $30+3=33$ couplages en moins (chaque bit d'écart avec les 10 bits de conception, et entre eux). Problème mis à l'échelle : 33 contre 30 qubits, et $3\times30+3=93$ couplages en moins.

\section*{Chapitre~\ref{ch:echelle} — Problème mis à l'échelle}
\textbf{Ex.~8.1.} (a) Près de $x=0$, $C(x)\approx\sqrt x$ et $S(0)=A_0$ (seul le premier polynôme de Bernstein vaut 1 en 0), donc $y\approx A_0\sqrt x$ et $y^2\approx A_0^2x$. Le cercle de rayon $R$ tangent en 0 donne $y^2\approx2Rx$ : $R=A_0^2/2$. (b) $C(1)=0$, donc $y(1)=0$ : le bord de fuite est fermé. Pour un bord de fuite épais, on ajoute à chaque face un terme $x\,\Delta y_{\rm BF}/2$.

\textbf{Ex.~8.2.} $8^{10}\approx1{,}07\cdot10^9$ secondes, soit environ 34 ans de simulations : hors de portée. $2^{33}\approx8{,}6\cdot10^9$ énergies à $10^8$ par seconde : environ 90 secondes, ce qui est faisable pour un QUBO. Mais la boucle en résout environ 300 par exécution (5 par itération), soit plus de 7 heures, et il faudrait stocker ou trier des milliards de valeurs. Ce qui rend le problème « grand », c'est la simulation de chaque profil ; le QUBO reste abordable bien plus longtemps. C'est encore le message du cours de calcul quantique : le coût est dans la simulation.

\section*{Chapitre~\ref{ch:resultats-v2} — Résultats V2}
\textbf{Ex.~7.1.} En médiane, $95-75=20$ simulations de plus (la médiane des écarts graine par graine vaut 22,5). Pour une graine donnée, les deux méthodes partent des \emph{mêmes} 10 profils initiaux : une graine « difficile » l'est pour les deux. Comparer graine par graine retire cette variabilité commune et ne garde que l'effet de la méthode. Le test apparié est donc plus sensible que la comparaison de deux médianes calculées sur des tirages qui varient beaucoup.

\textbf{Ex.~7.2.} $75\times(1{,}2-0{,}89)^2=75\times0{,}31^2\approx7{,}2$, soit environ sept fois l'étendue de l'objectif normalisé : aucun gain d'objectif ne peut compenser, et le profil n'est jamais proposé. Pour qu'il le soit, il faudrait que la pénalité reste de l'ordre d'un dixième de l'étendue, soit $1{,}2-\widehat{C_{L,\max}}\lesssim\sqrt{0{,}1/75}\approx0{,}04$ : une erreur de prédiction de moins de 0,04 sur une grandeur irrégulière. Autre solution : relâcher la contrainte là où le modèle est incertain (probabilité de faisabilité), ou utiliser un modèle plus indulgent pour la contrainte seule.

\section*{Chapitre~\ref{ch:resultats-echelle} — Résultats à l'échelle}
\textbf{Ex.~9.1.} (a) Sous l'hypothèse nulle, chaque différence a autant de chances d'être positive que négative, et les $2^8=256$ configurations de signes sont équiprobables. La plus extrême est celle où les 8 différences ont le même signe (statistique de Wilcoxon nulle) ; en bilatéral, on compte les deux sens : $p=2/256\approx0{,}0078$. Aucune configuration ne peut donner moins avec 8 graines. (b) Test des signes : $P(X\ge6)=\big(\binom86+\binom87+\binom88\big)/256=(28+8+1)/256\approx0{,}145$, donc $p\approx0{,}29$ en bilatéral. (c) Avec 20 graines et 15 gagnées : $P(X\ge15)=(15\,504+4\,845+1\,140+190+20+1)/2^{20}\approx0{,}021$, soit $p\approx0{,}041<0{,}05$. Avec 16 graines et 12 gagnées, $p\approx0{,}077$ : il faut une vingtaine de graines.

\textbf{Ex.~9.2.} (a) $300\times(1-0{,}59-0{,}14)=81$ simulations admissibles par graine. (b) 36 profils faisables pour 66 coefficients : le système est sous-déterminé. Avec $\lambda=10^{-6}$, le modèle en $z$ prend la solution de norme minimale (cours de calcul quantique, chapitre~8), qui passe par les données mais dit n'importe quoi ailleurs ; on est avant le pic « juste déterminé » du chapitre~\ref{ch:precision}, où la corrélation de rang sur $C_D$ reste vers 0,37. (c) La régularisation (pénalité $L_2$, a priori bayésien) ramène vers zéro les coefficients que les données ne déterminent pas. Dans une région inexplorée, les prédictions restent donc modérées au lieu d'extrapoler vers de très bonnes valeurs : le modèle propose moins de profils extrêmes, et moins de géométries invalides.

\section*{Chapitre~\ref{ch:bilan} — Bilan}
\textbf{Ex.~11.1.} Une réponse possible : « On approche la traînée par un polynôme de degré 2 des paramètres de forme, et l'on choisit le prochain profil en minimisant ce polynôme. Avec 4 paramètres, le polynôme a 15 coefficients, et après 75 simulations on en a cinq fois plus que de coefficients : il est bien déterminé, et le fait de savoir que les niveaux d'un paramètre se suivent est une information gratuite. Avec 10 paramètres, il a 66 coefficients, et l'on n'a guère plus de profils faisables que de coefficients : le polynôme est mal déterminé et peut extrapoler n'importe où. Les autres modèles ajoutent une régularisation, qui les rend prudents quand les données manquent. Ils rattrapent donc leur retard quand la dimension augmente. »

\textbf{Ex.~11.2.} L'objection est juste : avec une population de 16 et 300 simulations, NSGA-II fait moins de 20 générations, ce qui est peu pour un algorithme évolutionnaire. Battre NSGA-II montre que le modèle apporte de l'information, pas que la boucle bat les meilleures méthodes classiques. Il faut ajouter un concurrent conçu pour les petits budgets : une optimisation bayésienne, adaptée à la grande dimension (TuRBO, SAASBO). Sur la V2, le lot L1 a déjà montré qu'une optimisation bayésienne fait jeu égal avec la boucle. L'argument honnête est donc : « la boucle QUBO est au niveau des bonnes méthodes à modèle, et sa forme la rend compatible avec une machine quantique ».
